

\hypertarget{fixture-apa7-parity}{%
\section{Fixture: apa7 parity}\label{fixture-apa7-parity}}

The foundational treatment of pragmatic markers in contact situations
came from Smith (2010), whose framework has shaped subsequent work.
Recent quantitative evidence reinforced the general claim (Smith, 2010,
p. 22). Classroom-discourse data produced complementary findings (Jones
\& Cruz, 2015), while the urban multilingual synthesis broadened the
empirical base (Able et al., 2019).

A consolidated view across the three main lines of evidence is now
available (Able et al., 2019; Jones \& Cruz, 2015; Smith, 2010).

Institutional reporting on classroom multilingualism adds context from
outside the academic literature (Modern Language Association, 2021).
Undated industry commentary touches on similar themes (Perennial, n.d.).

\hypertarget{references}{%
\section*{References}\label{references}}
\addcontentsline{toc}{section}{References}

\hypertarget{refs}{}
\begin{CSLReferences}{1}{0}
\leavevmode\vadjust pre{\hypertarget{ref-able-baker-chen-2019}{}}%
Able, A., Baker, B., \& Chen, C. (2019). Multilingual Practices in Urban
Contexts. \emph{Language in Society}, \emph{48}(1), 55--80.

\leavevmode\vadjust pre{\hypertarget{ref-jones-cruz-2015}{}}%
Jones, R., \& Cruz, E. (2015). Code-switching in Classroom Discourse.
\emph{Applied Linguistics}, \emph{36}(2), 84--102.

\leavevmode\vadjust pre{\hypertarget{ref-mla-2021}{}}%
Modern Language Association. (2021). \emph{How do I cite a source with
multiple authors?} MLA Style Center.
\url{https://style.mla.org/citing-multiple-authors/}

\leavevmode\vadjust pre{\hypertarget{ref-perennial-nd}{}}%
Perennial, P. (n.d.). \emph{Undated Commentary}. Perennial Blog.
Retrieved \url{https://example.com/undated}

\leavevmode\vadjust pre{\hypertarget{ref-smith-2010}{}}%
Smith, J. A. (2010). Pragmatic Markers in Contact Situations.
\emph{Journal of Pragmatics}, \emph{42}(3), 12--34.
\url{https://doi.org/10.1016/j.pragma.2010.01.005}

\end{CSLReferences}

